\chapter*{Abstract}
\addcontentsline{toc}{chapter}{Abstract}

Exotic spin-dependent interactions mediated by light, weakly coupled bosons
(axions, axion-like particles, dark photons) are a well-motivated
low-energy signature of physics beyond the Standard Model. Two languages
are used to describe them: the Standard Model Extension (SME), a
relativistic effective field theory of Lorentz and CPT violation, and the
Dobrescu--Mocioiu (DM) classification, a complete basis of sixteen
non-relativistic potentials ($V_1$--$V_{16}$) in which experiments actually
report their bounds. No complete, systematic dictionary previously
connected the two, and no framework existed for using such a dictionary to
compare matter-sector and antimatter-sector constraints as a quantitative
test of CPT symmetry. This thesis develops both.

The Foldy--Wouthuysen transformation is applied to the SME fermion-sector
coefficients $b_\mu$, $d_{\mu\nu}$, $g_{\lambda\mu\nu}$ and $H_{\mu\nu}$ to
derive the non-relativistic Hamiltonian each generates and match it onto the
corresponding DM potential(s). These four are shown to be the complete set
that generates spin-dependent non-relativistic potentials: $a_\mu$,
$c_{\mu\nu}$ and $e_\mu$ appear only in spin-independent terms and $f_\mu$
not at all through third order in $1/m$, so none has an image in the DM
basis. All four
derivations are obtained in full and agree with established results,
including the momentum-dependent $d_{ij}$ and $d_{00}$ components, whose
apparent sign discrepancy against the literature is shown to arise from a
momentum-index convention in the reference result rather than from the
derivation. A dimensionless CPT asymmetry parameter, $A_\alpha =
(g_\alpha^{f}-g_\alpha^{\bar f})/(g_\alpha^{f}+g_\alpha^{\bar f})$, is
defined and, using a purpose-built pipeline (SPINDEP), evaluated against a
compiled database of 283 experimental constraint datasets, 247 of which
enter the analysis, spanning
nitrogen-vacancy-centre magnetometry, torsion balances, atomic
co-magnetometers, positronium spectroscopy, and antimatter experiments,
standardised to a common interaction-range axis.

Only 22 of the 247 analysed datasets involve an antimatter sector, from
which fifteen matter--antimatter pairs could be statistically compared. Every pair shows
a significant asymmetry after correcting for autocorrelation among
interpolated points, yet an analytic argument, confirmed empirically across
all fifteen pairs, shows that $A_\alpha$ built from one-sided experimental
upper bounds cannot by itself distinguish genuine CPT violation from an
ordinary sensitivity gap between the matter- and antimatter-sector
experiments being compared. A gap analysis by fermion sector and by DM
potential identifies the $e$--$N$, $n$--$N$, $p$--$N$, $e$--$n$, $n$--$p$,
and $n$--$n$ sectors,
and the $V_{1a}$, $V_{9+10}$, and $V_{12+13}$ potentials, as the largest
opportunities for new antimatter-sector measurements. The thesis concludes
that the antimatter-sector coverage gap, not statistical power, is the
binding constraint on testing CPT symmetry through exotic spin-dependent
interactions, and offers the mapping table, the database, and the analysis
pipeline as an extensible resource for closing it.

\textbf{Keywords:} Standard Model Extension, Dobrescu--Mocioiu potentials,
Foldy--Wouthuysen transformation, CPT symmetry, exotic spin-dependent
interactions, matter--antimatter asymmetry.
